\section{Recommendations for Multi-Agent Workflows}

The foundations chapter described role specialization as trading the simplicity of a single continuous loop for narrower, more focused reasoning at each stage, coordinated through structured hand-off of artifacts rather than open-ended conversation --- and it flagged a specific cost that comes with this trade: an error introduced early in the pipeline can silently propagate through every later stage rather than being caught immediately, particularly if a later agent has no way of seeing the original task and can only work from what the previous agent handed it. Part 2's Architect--human--Implementer--Reviewer sequence is a concrete instance of this pattern, and the way it was structured is worth drawing recommendations from directly, since it addresses the propagation risk rather than simply accepting it.

The role boundaries themselves were kept deliberately narrow and non-overlapping: the Architect Agent is explicitly barred from writing implementation code, limited to producing plans, interface definitions, and file-structure proposals, while the Implementer is restricted to creating and modifying exactly the files the approved plan lists, and is required to stop and request an updated plan rather than resolve any gap the plan never addressed itself. \textbf{A recommendation that follows from this is to define each agent role by what it is not permitted to do as much as by what it is responsible for, since a role with authority to quietly expand beyond its planning or implementation boundary reintroduces the single-agent problem of one actor carrying every concern at once, inside what looks like a multi-agent structure.}

The propagation risk the foundations chapter describes is addressed most directly by where the human gate sits relative to the hand-off between roles, rather than by anything internal to how the agents communicate. The plan's Design Review is required to contain five specific sub-sections --- architecture rationale, state ownership, service ownership, testing strategy, and scalability concerns --- and a plan missing any of them is treated as invalid and cannot be approved, which means the human reviewing the hand-off from Architect to Implementer has a fixed, comparable structure to check each time rather than an open-ended document to interpret freely. \textbf{This suggests that a multi-agent workflow's hand-off points, not just its individual roles, need an explicit, enforced structure to review against; without one, a human gate can be nominally present while still allowing an incomplete or vague plan through, at exactly the point where an error would otherwise begin propagating into every stage built on top of it.}

The workflow also builds in a mechanism for what happens when an agent's scope proves too narrow for something it encounters --- a case the foundations chapter's error-propagation concern does not by itself resolve. Rather than having the Implementer either overstep its boundary to fix an unplanned issue or silently ignore it, the workflow requires the issue to be documented as its own candidate task, to be picked up through the same plan-review-implement sequence rather than folded into whatever is currently in progress. \textbf{The practice this generalizes is to give every agent role an explicit escape valve for work that falls outside its defined scope --- a required documentation step rather than an implicit choice between overstepping the boundary or leaving the issue unaddressed --- since narrow roles are only safe to enforce strictly if there is a defined place for what they are not allowed to handle themselves.}

